\chapter{Типизированный мост энергии часов к скорости полного шума}
\label{ch:v8-typed-clock-energy-to-noise-rate-anchor-gate}

Автономные часы уже обеспечивают локальный непрерывный предел, но их
частота и скорость диссипации допускают общий пересчёт. Проверим, какую
размерную формулу действительно даёт столкновительная модель и какие данные
ещё должен выбрать родитель.

\section{Размерная формула}

Пусть энергия часов равна $E_C$, длительность такта
\begin{equation}
 \tau_C=\frac{\hbar}{E_C},
\end{equation}
а амплитуда энергии взаимодействия одного столкновения имеет вид
\begin{equation}
 E_{\rm int}=\chi E_C,
\end{equation}
где $\chi>0$ безразмерна. Квадратичный слабый предел даёт
\begin{equation}
 \Gamma=\frac{E_{\rm int}^2\tau_C}{\hbar^2}
 =\chi^2\frac{E_C}{\hbar}.
 \label{eq:v8-typed-clock-noise-rate}
\end{equation}
Для частоты $\Omega=E_C/\hbar$ получаем точное отношение
\begin{equation}
 \boxed{\frac{\Gamma}{\Omega}=\chi^2.}
 \label{eq:v8-clock-rate-relative-calibration}
\end{equation}

\begin{theorem}[Условная относительная калибровка]
Если энергия часов $E_C$ и безразмерное отношение
$\chi=E_{\rm int}/E_C$ заданы одним типизированным микроскопическим
родителем, то часы фиксируют скорость полного шумового процесса формулой
$\Gamma=\chi^2E_C/\hbar$.
\end{theorem}

\section{Что текущий родитель не выбирает}

Нынешняя конструкция фиксирует форму 42 операторов скачка, но допускает
любое положительное умножение взаимодействия. Часовой гамильтониан введён
отдельно и также допускает изменение энергетического масштаба. Поэтому два
значения $\chi_1\ne\chi_2$ сохраняют тот же носитель часов, тот же шумовой
кадр и все симметрии, но дают разные отношения
\begin{equation}
 \frac{\Gamma_1}{\Omega}=\chi_1^2,
 \qquad
 \frac{\Gamma_2}{\Omega}=\chi_2^2.
\end{equation}
Приравнивание $E_{\rm int}=E_C$ означало бы дополнительное условие
$\chi=1$, а не следствие уже построенного действия.

Старые спектральные щели, массы и параметр обрезания не исправляют этот
дефект автоматически. В текущей цепи нет морфизма, который переводил бы их
в коэффициент часового гамильтониана и одновременно в амплитуду
столкновения. Без такого морфизма они остаются внешними кандидатами на
калибровку.

\begin{nogocriterion}[Отсутствие размерного якоря в текущем родителе]
Формула \eqref{eq:v8-typed-clock-noise-rate} типизирована, но её два входа
$E_C$ и $\chi$ не выбраны. Без нового смешанного члена или независимо
измеренного энергетического масштаба нельзя объявлять ни $\Gamma$, ни
$\tau_C$ физически предсказанными. Безразмерная щель генератора сама по себе
энергией не является.
\end{nogocriterion}

\begin{protocol}
\begin{enumerate}
 \item длительность такта: \textbf{$\tau_C=\hbar/E_C$};
 \item энергия взаимодействия: \textbf{$E_{\rm int}=\chi E_C$};
 \item скорость шума: \textbf{$\Gamma=\chi^2E_C/\hbar$};
 \item относительная калибровка: \textbf{$\Gamma/\Omega=\chi^2$};
 \item значение $E_C$: \textbf{не выведено};
 \item значение $\chi$: \textbf{не выведено};
 \item условие $E_{\rm int}=E_C$: \textbf{не допускается без родителя};
 \item старая щель как энергия: \textbf{не типизирована};
 \item абсолютная скорость $\Gamma$: \textbf{не выведена};
 \item следующий гейт: \textbf{аудит кандидатов на энергетический якорь часов}.
\end{enumerate}
\end{protocol}

Машинный сертификат сохранён в
\path{s2t_v8_typed_clock_energy_to_noise_rate_anchor_gate_results.json}.

\begin{thebibliography}{9}
\bibitem{v8-typed-clock-attal}
S.~Attal, Y.~Pautrat,
\emph{From repeated to continuous quantum interactions},
Ann. Henri Poincar\'e \textbf{7} (2006), 59--104;
arXiv:math-ph/0311002.
\bibitem{v8-typed-clock-bruneau}
L.~Bruneau, A.~Joye, M.~Merkli,
\emph{Repeated interactions in open quantum systems},
J. Math. Phys. \textbf{55} (2014), 075204; arXiv:1305.2472.
\end{thebibliography}